\chapter{Theoretical Background}
\section{Central Bank Communication}
% opísať, čo banka komunikuje s verejnosťou, čo sa vie dozvedieť verejnosť
% že existuje viacero druhov konferencií a že v bakalárke budeme pracovať s MP statement

\subsection{Monetary Policy Statement}
% rozvrh dňa, ako často sa deje a tak
% že sa delí na 2 časti - intro a Q&A
% intro je vysoko štrukturizovaný zatiaľ čo Q&A je take a take
% proste také vysvetlenie čo to je vôbec tento statement
% +obrázok ako príklad, že ako vyzerá taký statement

\subsection{Other Communication} %môžeš zmeniť nadpis podľa kontextu
% ďalšie druhy komunikácii (mesačný buletin...)

\section{Natural Language Processing (NLP) in Finance}
% základná teória k tomu + jemný úvod k ďalším podsekciám

\subsection{Sentiment Analysis}
% čo to je a načo sa používa, možno spomenúť finančné texty
% tu nezabudnúť o pozitiv vs neutral vs negatív

\subsection{BERT, FinBERT, etc.}
% rôzne modely, ktoré sa používajú pre analýzu finančných textov (medzi bankovými)
% doterajšie zistenia a výsledky v analýze textov z ECB / FED
% nezabudnúť, že skôr sa hovorí hawish vs dowkish tón
% ideálne iba tie, ktoré vyskúšam

\section{Time Series Analysis}
% v skratke čo to je
% v praktickej časti:
% - mám text -> sentimenty (rozdelený text na viac častí a z neho sentimenty)
% - mám spracované time-series -> spárujem to so sentimentom
% - potom trénujem model -> aby predvídal nejaké rasty / poklesy a tak mal výnos

\section{Portfolio Optimization}
% efektívne riadenie portfólia